% SGA 3，第 XI 讲；根据英文重建源翻译。
% 范围：命题 1.9 及说明、命题 1.10 及证明、推论 1.11 及证明；止于第 2 节之前。
% 权威来源：Polo--Gille Expo11.pdf，组件局部页 6--7；合卷页 728--729；印刷／原文页边页 135--137。
% Jacob Reinhold 的 CC BY 4.0 Markdown 谱系仅用于措辞比较；文字、公式、编号、注释及图均以 Polo--Gille PDF 为准。
% 精确续接／终点：第 2 节，Expo11 局部页 7／合卷页 729／印刷页 137。

\medskip\noindent\SGARefTarget{sga3:xi:proposition:1:9}\textbf{命题 1.9。} --- \textit{设 \(f:X\to S\) 为预概形的平滑态射。则 \(\Omega^1_{X/S}\) 是 \(X\) 上有限型的局部自由模，且其在点 \(x\in X\) 处的秩等于纤维 \(X_s\) 在 \(x\) 的某个邻域内的维数，其中 \(s=f(x)\)。}

此维数称为\emph{\(X\) 在 \(x\) 处相对于 \(S\) 的维数}。注意（当 \(f\) 在 \(x\) 处平滑时），它为零当且仅当 \(f\) 在 \(x\) 处平展。也可以用 \(X\) 所表示的函子 \(F\) 直接计算此维数，方法如下。设 \(\xi\) 为 \(F\) 的一个取值于域 \(k\)、“位于 \(x\) 处”的点，即像为 \(x\) 的态射 \[
  \operatorname{Spec}(k)\longrightarrow F=X
\]。考虑 \(k\) 上的对偶数代数 \[
  D(k)=k[t]/(t^2)
\]，将其视为 \(S\)-预概形，并考虑由增广 \(D(k)\to k\) 诱导的映射 \SGARefTarget{sga3:xi:display:1:9:dual-number-map} \[
  F(D(k))\longrightarrow F(k)
\]。最后，设 \(F(D(k),\xi)\) 为此映射下 \(\xi\in F(k)\) 的逆像。此集合自然具有 \(k\) 上的向量空间结构。事实上，它是 \[
  \Omega^1_{X/S}\otimes_{\mathscr O_{X,x}}k,
\] 的对偶向量空间，其维数就是 \(X\) 在 \(x\) 处相对于 \(S\) 的维数。

为明确描述 \(F(D(k),\xi)\) 上的向量空间运算，宜如第二讲那样，更一般地对 \(k\) 上的每个向量空间 \(V\) 引入代数 \[
  D_k(V)=k+V,
\]，其中 \(V\) 为平方为零的理想，并将 \[
  F(D_k(V))\longrightarrow F(k),
\] 下 \(\xi\) 的逆像 \(F(D_k(V),\xi)\) 视为 \(V\) 的取值于 \((\mathbf{Ens})\) 的协变函子。只须此函子与两个因子的乘积交换。这意味着 \(F\) 将某些非常特殊的融合和变为纤维积；参照第二讲。当 \(F\) 可表时，此条件总成立。因此，各集合 \(F(D_k(V),\xi)\)，特别是 \(F(D(k),\xi)\)，具有 \(k\) 上的向量空间结构。所以，在远比 \(F\) 可表更一般的条件下，也可以定义 \(F\) 在“点” \(\xi\) 处相对于 \(X\) 的维数。

本讲中，某些将明确描述的函子可由 \(S\) 上平滑的预概形表示这一事实，主要通过 \SGARefLink{sga3:xi:proposition:1:10}{下述结果} 使用。它将成为我们的技术手段，使我们从 Noether 概形 \(S\) 的点 \(s\) 的局部环 \(A\) 的完备化上的构造，过渡到 \(A\) 上、在 \(A\) 上平展的局部环 \(A'\) 上；特别地，得以从 \(s\) 过渡到邻近点。

\medskip\noindent\SGARefTarget{sga3:xi:proposition:1:10}\textbf{命题 1.10。} --- \textit{设 \(f:X\to S\) 为预概形的平滑态射，\(s\in S\)，\(x\in X\) 位于 \(s\) 上方，且 \(\kappa(x)\) 为 \(\kappa(s)\) 的有限可分扩张。则存在经过 \(x\)、在 \(S\) 上平展的子概形 \(S'\subset X\)。\footnote{法文再版印作“\(S\) 的一个子概形 \(S'\)，在 \(S\) 上平展，经过 \(x\)”。由于 \(x\in X\)，且下面的等价表述与构造均使用 \(S'\to X\) 和 \(\mathscr O_X\) 的截面，所意指的应是 \(S'\subset X\)。}因此，可以找到平展态射 \(S'\to S\)、\(s\) 上方的点 \(s'\in S'\)，其剩余域扩张为 \(\kappa(x)/\kappa(s)\)，以及将 \(s'\) 映到 \(x\) 的 \(S\)-态射 \(S'\to X\)。}

为构造 \(S'\)，只需在 \(x\) 的某个邻域上选取 \(\mathscr O_X\) 的截面 \(g_1,\ldots,g_n\)，使其在 \(x\) 处的像构成纤维 \(X_s\) 在 \(x\) 处的局部环的正则参数系。于是由各 \(g_i\) 定义的子概形 \(S'\) 在 \(x\) 处于 \(S\) 上平展，缩小 \(S'\) 后，它在 \(S\) 上平展。

我们将在 \(\kappa(x)=\kappa(s)\) 时使用 \SGARefLink{sga3:xi:proposition:1:10}{1.10}，即 \(x\) 在 \(\kappa(s)\) 上有理，或等价地，\(x\) 可视为 \(X_s\) 在 \(\operatorname{Spec}\kappa(s)\) 上的截面。此时，\SGARefLink{sga3:xi:proposition:1:10}{1.10} 具有在作平展基扩张后延拓截面的定理的性质。在以下特殊情形中，它具有特别简单的形式。

\medskip\noindent\SGARefTarget{sga3:xi:corollary:1:11}\textbf{推论 1.11。} --- \textit{在 \SGARefLink{sga3:xi:proposition:1:10}{1.10} 的假设下，假设 \(S\) 为 Hensel 局部环的谱，且 \(\kappa(x)=\kappa(s)\)。则存在 \(X\) 在 \(S\) 上经过 \(x\) 的截面。若 \(X\) 实际上在 \(x\) 处于 \(S\) 上平展，则此截面唯一。}

事实上，由于 \(S\) 为 Hensel 概形，在 \SGARefLink{sga3:xi:proposition:1:10}{1.10} 的假设下，预概形 \(S'\) 包含一个在 \(S\) 上有限、且在 \(s\) 处的纤维仅为 \(x\) 的开子概形。由于此开子概形在 \(S\) 上平展，它同构于 \(S\)，断言得证。注意，当 \(S\) 为完备局部环的谱时，\SGARefLink{sga3:xi:proposition:1:10}{1.10} 或 \SGARefLink{sga3:xi:corollary:1:11}{1.11} 大致等价于经典的“Hensel 引理”，有时也以此名指称该结果。


